% Generated draft: requires mathematical and source review.
\documentclass[11pt]{article}
\usepackage[margin=25mm]{geometry}
\usepackage{fontspec}
\setmainfont{Latin Modern Roman}
\usepackage{amsmath,amssymb,mathtools,unicode-math}
\setmathfont{Latin Modern Math}
\usepackage{enumitem,needspace,xurl,hyperref}
\setlength{\parindent}{0pt}
\setlength{\parskip}{5pt}
\setlength{\emergencystretch}{3em}
\newcommand{\N}{\mathbb{N}}
\newcommand{\Z}{\mathbb{Z}}
\newcommand{\Q}{\mathbb{Q}}
\newcommand{\R}{\mathbb{R}}
\newcommand{\C}{\mathbb{C}}
\newcommand{\F}{\mathbb{F}}
\newcommand{\A}{\mathbb{A}}
\newcommand{\PP}{\mathbb{P}}
\newcommand{\E}{\mathbb{E}}
\newcommand{\eps}{\varepsilon}
\newcommand{\dd}{\,\mathrm{d}}
\begin{document}
\section*{UnsolvedMath ID: 30004673}
\textbf{Source problem code:} catalogue entry
\subsection{Definitions and mathematical statement}
The supplied record does not define Poisson genericity or specify the space, representation, or probability measure for its reals. Consequently, the following is a schematic formulation whose missing data require mathematical review.

Let \(X\) denote the intended space of represented reals, and let \(P\subseteq X\) denote the intended class of Poisson-generic reals. Here \(P\) is a placeholder for the missing definition. Once an effective representation of \(X\) is fixed, write \(\operatorname{Comp}(X)\) for its computable elements. Once the probability measure and effective test conventions are fixed, write \(\operatorname{MLR}(X)\) for the corresponding Martin--L\"of random elements.

The source asks the following four questions:
\begin{enumerate}
\item Does a computable Poisson-generic real exist; equivalently, is
\[
P\cap\operatorname{Comp}(X)\ne\varnothing?
\]
\item Is every Martin--L\"of random real Poisson generic; equivalently, does
\[
\operatorname{MLR}(X)\subseteq P
\]
hold?
\item Find a necessary and sufficient characterization of membership in \(P\) in terms of Kolmogorov complexity. The source does not specify the complexity variant, the representation of finite data extracted from a real, or the permissible form of the characterization.
\item Is \(P\) \(\Pi^0_3\)-complete? The source does not specify whether this denotes an effective arithmetical classification or a Borel classification, nor the reductions under which completeness is intended.
\end{enumerate}

The definition of Poisson genericity, its parameters and quantifiers, and the conventions above must be supplied before these questions constitute a fully specified mathematical statement. The supplied dated curation labels the problem partially solved: it reports affirmative results for the first two questions, settlement of the completeness question without stating its answer, and an unresolved complexity-characterization question in the requested form. These are supplied literature assessments, not independently verified mathematical conclusions.
\subsection{Short English statement}
The problem asks whether Poisson genericity admits a computable example, contains every Martin--L\"of random real, has a Kolmogorov-complexity characterization, and defines a \(\Pi^0_3\)-complete class. The supplied record omits the definition and conventions needed for a precise formulation. Its dated curation reports partial resolution, with the requested complexity characterization still unresolved in this form.
\subsection{Sources}
\begin{itemize}
\item Local UnsolvedMath registry: \texttt{data/unsolvedmath/problems.json}, permanent ID \(30004673\), supplied source record titled ``Computability and Complexity Characterizations of Poisson Genericity,'' source problem code \texttt{catalogue entry}. The record belongs to the Oberwolfach Reports Open Problems collection and records proposed year 2021; no individual proposer is supplied.
\item Supplied report citation: \emph{Computability Theory (2022)}, Oberwolfach Reports, pp.~1149--1190, DOI \texttt{10.4171/owr/2021/21}. \url{https://doi.org/10.4171/owr/2021/21}. The background describes this collection as extracted from workshop reports and their open-problem discussions.
\item \textrm{\'Alvarez}, Becher and Mereb, ``Poisson generic sequences.'' \url{https://doi.org/10.1093/imrn/rnac234}.
\item Becher, Jackson, Kwietniak and Mance, supplied reference label ``descriptive complexity.'' \url{https://arxiv.org/abs/2305.10529}.
\item Oberwolfach Report 21/2021. \url{https://ems.press/content/serial-article-files/46899}.
\item Literature assessment date: \texttt{2026-08-22}. This is supplied dated curation derived from the OWR curation, not an independently verified assessment. It reports that computable Poisson-generic reals exist, that all Martin--L\"of random reals are Poisson generic, that the \(\Pi^0_3\)-completeness question was settled, and that the requested Kolmogorov-complexity characterization remains unresolved in this form. The supplied background explicitly calls for independent verification of time-sensitive status and cited claims.
\item The supplied statement status is \texttt{corrected\_verified}; its accompanying qualification says that only formula delimiters, notation, grammar, and typography were repaired, with mathematical content and scope preserved. This label does not provide the missing definitions or constitute mathematical review of this draft.
\end{itemize}
\end{document}
